% Corrigé intégré automatiquement pour la banque Exercices.
% Source : A_Integrer/DDS_04/079_Geometrie_Verin/079_Geometrie_Verin.tex
\begin{corrigebox}[Corrigé]
\CorrigeQuestion{1}
La résolution consiste à identifier les données et l'inconnue, choisir la loi du cours adaptée, écrire l'équation symbolique avant toute application numérique, puis vérifier l'homogénéité, le signe et l'ordre de grandeur du résultat. Les hypothèses utilisées doivent être explicitement contrôlées à la fin.
\CorrigeQuestion{2}
On définit les variables, le pas de calcul et les conditions initiales, puis on traduit l'équation par une relation de récurrence. L'algorithme est vérifié sur les premiers pas et la stabilité numérique est contrôlée en comparant le pas aux constantes de temps du système.
\CorrigeQuestion{3}
La résolution consiste à identifier les données et l'inconnue, choisir la loi du cours adaptée, écrire l'équation symbolique avant toute application numérique, puis vérifier l'homogénéité, le signe et l'ordre de grandeur du résultat. Les hypothèses utilisées doivent être explicitement contrôlées à la fin.
\CorrigeQuestion{4}
On définit les variables, le pas de calcul et les conditions initiales, puis on traduit l'équation par une relation de récurrence. L'algorithme est vérifié sur les premiers pas et la stabilité numérique est contrôlée en comparant le pas aux constantes de temps du système.
\end{corrigebox}
